\section{Related Work}\label{sec:related}

\paragraph{Long-term memory for agents.} Recent work treats agent memory as an
architecture rather than a vector store. ZenBrain proposes a seven-layer,
neuroscience-derived memory architecture with sleep-phase consolidation and
explicit ablations of each mechanism~\cite{zenbrain2026}. FSFM argues that
selective forgetting deserves first-class treatment and gives a taxonomy of
passive decay, active deletion, safety-triggered forgetting, and adaptive
reinforcement~\cite{fsfm2026}. A human-inspired architecture combines
sleep-phase consolidation, interference-based forgetting, and multi-cue
retrieval, reporting retention precision of 97.2\% at a 58\% store
reduction~\cite{himemory2026}. PersMem internalizes personality into a
dual-pathway memory pipeline so that memory processing becomes
personality-dependent~\cite{persmem2026}. Mem0 extracts, consolidates and
retrieves salient conversational facts at production scale, reporting large
latency and token savings over full-context prompting~\cite{mem0_2025}. These
systems optimize what is remembered and how well it is retrieved. Our concern is
different and orthogonal: what is \emph{allowed} to enter the prompt, under which
priority, and what happens when the budget cannot hold everything. We use Mem0 as
an external memory baseline for that question
(Section~\ref{sec:exp:baselines}).

\paragraph{Activation and decay.} Ranking memories over long horizons is usually
framed as a decay or activation problem. The ACT-R base-level activation folds
frequency and recency into a single logarithmic quantity, motivated by the
statistics of practice and need probability~\cite{anderson1991,anderson2007}.
We treat it as one of two channels rather than as a replacement for decay,
because an outright replacement degrades a frozen ranking baseline in our own
setting, as reported in Section~4. This is a concrete example of a general
point: borrowed mechanisms rarely transplant without a regression check against
the instrument that existed before the borrowing.

\paragraph{Affect and persona.} Sentipolis couples a continuous
Pleasure--Arousal--Dominance representation with dual-speed emotion dynamics
and emotion--memory coupling for social simulation~\cite{sentipolis2026}. Its
own report is instructive: gains are model-dependent, and emotion-awareness can
mildly reduce adherence to social norms. PAD itself dates back to environmental
psychology~\cite{mehrabian1974}. We adopt PAD as a representation and, unlike
simulation-oriented systems, confine affect to the style domain: an invariant
forbids affect from entering factual judgement, and the composer enforces the
separation structurally.

\paragraph{Safety engineering.} STPA analyses unsafe control actions in
safety-critical systems~\cite{leveson2011,leveson2018}. We adapt the notion of
an unsafe control action from its original, offline analytical setting to a
runtime gate: eight control actions are screened after composition and before
the model call, with a single fail-closed locus. The transfer is deliberate but
partial---STPA is a hazard-analysis method, whereas our gate is a predicate
evaluated on every turn---and we do not claim that our gate is a substitute for
a full STPA of the host.

\paragraph{Cybernetics and requisite variety.} Ashby's law of requisite variety
states that a regulator must command at least as much variety as the
disturbances it faces~\cite{ashby1956}; Beer's Viable System Model is one
influential reading of layered regulation~\cite{beer1972}. We borrow the
engineering idea that when a controller cannot represent a conflict, it should
shrink its own action space instead of acting anyway, and implement this as
dimensionality reduction over composition domains. We use VSM only as an
interpretive lens on an existing priority layering: it is borrowed, not
implemented, has no code locus, parameter, or dependency in our implementation,
and we do not claim System~1--5 machinery. Anti-windup for integral controllers
is standard control practice~\cite{astrom2006} and is adopted directly.

\paragraph{Evaluation and staged change.} Optimizing a single proxy invites
Goodhart's law~\cite{goodhart1975}, so our learning loop pairs every primary
metric with an adversarial counterpart, treats user corrections as the final
judge, and freezes the evaluation instrument itself. Change management follows
a conventional staged-rollout pattern (small share, guardrails, automatic
rollback), with the additional constraint that data insufficient to verify a
guardrail blocks release instead of skipping it.

\paragraph{Scope.} The mechanisms discussed above are the whole of the paper's
scope. Limits of that scope, including components that exist in the project but
are deliberately not claimed here, are stated in Section~7.
